%Au 1\up{er} janvier 2025, une réserve d’eau naturelle contient \np[m^3]{3000} d'eau.
%
%Une étude indique que la réserve perd chaque année 10~\% du volume d'eau qu'elle avait en début d'année.
%
%On modélise l'évolution du volume d'eau de cette réserve par une suite $\left(u_n\right)$, où $u_n$ désigne la quantité d'eau, exprimée en mètres cubes, contenue dans la réserve au 1\up{er} janvier de l'année $2025 + n$.
%
%On a donc $u_0 = \np{3000}$.

\begin{enumerate}
  \item 
%  Montrer que $u_1 = \np{2700}$.
Chaque année, le volume diminue de 10\%. Comme 10\% de \np{3000} c'est 300 (le dixième), le volume de la réserve sera égal à \mbox{$\np{3000}-300=\np{2700}$} au 1\up{er} janvier 2026, et donc \mbox{$u_1=\np{2700}$}.
  \item 
%  Pour tout entier naturel $n$, exprimer $u_{n+1}$ en fonction de $u_n$ .
Enlever 10\% d'une quantité revient à multiplier celle-ci par \mbox{$1-\dfrac{10}{100}=\np{0.9}$}.

Donc, pour tout \mbox{$n \in \N, u_{n+1}=\np{0.9}u_n$}.
  \item 
%  Quelle est la nature de la suite $(u_n)$ ? Préciser sa raison.
\mbox{$(u_n)$} est donc une suite géométrique (relation de récurrence multiplicative), et sa raison est \np{0.9}.
  \item 
%  Dans la feuille de calcul suivante on a calculé les premiers termes de la suite $(u_n)$. Le format des cellules de la ligne 2 permet d’afficher les nombres arrondis à l’unité.
%
%\begin{center}  
%\begin{minipage}{.6\linewidth}\renewcommand{\arraystretch}{1.15}
%\helvbx
%\begin{tblr}{
%	row{1} = {gray9,c,m},
%	row{2} = {c},
%	row{3} = {c},
%	column{1} = {gray9,c,m},
%	hlines,vlines,
%	}
%    & A & B & C & D & E & F & G & H & I  \\
%  1 & $n$ & 0 & 1 & 2 & 3 & 4 & 5 & 6 & 7  \\
%  2 & $u_n$ & 3000 & 2700 & 2430 & 2187 & 1968 & 1771 & 1594 & 1435 
%\end{tblr}
%\end{minipage}
%\end{center}
%
%Quelle formule a été saisie en cellule {\helvbx C2} afin d'obtenir, par glisser-coller vers
%la droite, les premiers termes de la suite $(u_n)$ ?
On a utilisé dans la cellule \textbf{C2} la relation de récurrence, pour écrire \texttt{=0,9*B2}.
  \item 
%  Pour assurer la sécurité de la lutte contre les incendies on estime que la réserve doit contenir au moins \np[m^3]{2000}.
%
%D’après le tableau précédent, à partir de quelle année la quantité d’eau dans la réserve deviendra insuffisante ?
D'après le tableau, la quantité d'eau devient inférieure à \np[~m^3]{2000} pour \mbox{$n=4$} : c'est donc à partir du 1\up{er} janvier 2029 que la quantité d’eau dans la réserve deviendra insuffisante.
\end{enumerate}  

\needspace{10\baselineskip}
% saut de page si moins de 10 lignes disponibles
